\section{Numerical algorithm, validation and cost}\label{sec:algorithm}

\subsection{Root evaluation without an explicit penalty}
Equation~\eqref{eq:bracket} provides a finite bracket
for any interior $s$ and a finite computed positive singular
spectrum. For numerical work, form
\[
\alpha_j=2\log\sigma_j-\log\delta_*,
\quad \log\delta_*=\operatorname{median}_j(2\log\sigma_j),
\]
rather than computing $\delta_j$ and their ratios directly.
Using the logistic function, evaluate
$G(x)=r^{-1}\sum_j\expit(x+\alpha_j)$.
When $s>1/2$, it is preferable to solve the
\emph{complementary equation}
\begin{equation}\label{eq:tail}
\frac1r\sum_{j=1}^r\expit(-x-\alpha_j)=1-s
\end{equation}
rather than subtracting two nearby numbers after evaluating $G(x)$.
Brent's bracketing method, bisection or safeguarded Newton iteration
can find $x_s$; the strict monotonicity and bracket in
Theorem~\ref{thm:root} eliminate ambiguous roots. These choices
do not imply a better condition number than log-$\lambda$
inversion.

\paragraph{Direct-$S$ reconstruction algorithm.}
Assume $D_d=U\Sigma V^\top$ has been computed and its positive
singular values pass the chosen numerical-rank check.
\begin{enumerate}
\item Accept $s\in[0,1]$, $y\in\R^n$, and either a fixed finite
$\mu$ or the explicit plug-in $\widehat\mu(y)$ of \eqref{eq:plug}.
\item If $s=0$, return $\widehat\tau=y$. If $s=1$,
set $w_j=1$ for every positive singular direction.
Otherwise compute the log-ratios $\alpha_j$ and
bracket $x$ using \eqref{eq:bracket}.
\item Solve $G(x)=s$ (or \eqref{eq:tail}) with
a bracket-preserving scalar algorithm. Set
$w_j=\expit(x+\alpha_j)$.
\item Form $c_\mu=V^\top y-\mu\Sigma^{-1}U^\top\one_r$
and return $\widehat\tau=y-V(w\odot c_\mu)$.
Optionally return the \emph{conditional} EDF
$\nu_{\rm cond}=n-rs$, the attained mean
$r^{-1}\sum_jw_j$ and a root residual.
For a plug-in mean, separately report
$\nu_{\rm plug}$ from \eqref{eq:edf-plug}.
\end{enumerate}
The first two endpoints use limiting weights, not matrix
factorization at an infinite or zero penalty. The mathematical
formula avoids forming the potentially large degree-$d$
particular polynomial $p_\mu$ in the interior.

\begin{corollary}[Residual-to-trend error certificate]\label{cor:certificate}
Suppose an approximate scalar solution $\widetilde x\in\R$
yields computed \emph{exact-arithmetic} weights
$\widetilde w_j=\expit(\widetilde x+\alpha_j)$
with scalar residual
$\varepsilon=G(\widetilde x)-s$.
Let $\widetilde\tau=y-V(\widetilde w\odot c_\mu)$.
Then
\[
\norm{\widetilde\tau-\widehat\tau_{s,\mu}}_2
\le r|\varepsilon|\norm{c_\mu}_2.
\]
\end{corollary}
\begin{proof}
Since $G$ is strictly increasing and maps $\R$ onto $(0,1)$,
$\widetilde w=w(\widetilde s)$ with
$\widetilde s=G(\widetilde x)=s+\varepsilon$.
Thus $\widetilde\tau=\widehat\tau_{\widetilde s,\mu}$;
apply \eqref{eq:lipschitz}.
\end{proof}

\begin{remark}[Scope of the certificate]
This is an \emph{algebraic root-residual} bound. A floating-point
implementation also incurs errors in its computed SVD, scalar
logistic sums and matrix products. The certificate alone
does not bound those errors, and it cannot establish that the
new approach is more accurate than computing the same
weights through a carefully stabilized log-$\lambda$
algorithm.
\end{remark}

\subsection{Arithmetic, storage and reuse}

Let $I$ denote the number of scalar residual evaluations.
For dense reduced SVD of an $r\times n$ matrix with $r\le n$,
standard algorithms take $O(nr^2)$ setup work and $O(nr)$
storage. Solving the scalar equation requires $O(Ir)$ work;
reconstructing one full $n$-vector after the SVD costs
$O(nr)$ for spectral projections and multiplication.
Calculating $\widehat\mu(y)=r^{-1}\one_r^\top D_dy$
requires up to $O(nd)$ arithmetic using a difference
stencil. No dense $n\times n$ inverse is formed.

If $K$ fixed smoothness targets are used on $R$ series of
the same length and difference order, the weights can in
principle be precomputed once and stored with $O(Kr)$
additional memory. The total major operations are
\begin{equation}\label{eq:cost}
O\bigl(nr^2+KIr+RKn r\bigr)
\end{equation}
for $RK$ complete reconstructions after shared preparation;
series-dependent projections may also be shared.
For only a single fit, structured banded factorization may
be asymptotically cheaper and should be treated as an
essential competitor \citep{weinert2007}.
A cached inverse-penalty implementation could cache
\emph{the identical weights}; thus \eqref{eq:cost}
is a description of the proposed execution architecture,
not a proof of superiority over a well-engineered baseline.

\begin{table*}[t]
\caption{Mathematical and computational status of three equivalent implementations. The final comparison must time each phase at identical $s,d,\mu$, precision and hardware, rather than infer a speedup from eliminating the name of $\lambda$.}
\label{tab:comparison}\centering
\begin{tabular}{@{}p{.23\textwidth}p{.25\textwidth}p{.27\textwidth}p{.14\textwidth}@{}}
\toprule
Execution path & Scalar operation & Reconstructing the trend & Current evidence\\
\midrule
Explicit $\Lambda_d(s)$ plus spectral PLS
& Invert $S_d(\lambda)$, retain $\lambda$ & Evaluate $w_j(\lambda)$ and shift-free SVD & Existing validated baseline\\
Implicit direct weights
& Solve $G(x)=s$, do not form $\lambda$ & Evaluate $w_j=\expit(x+\alpha_j)$ and shift-free SVD & Implemented; algebraically equivalent\\
Cached weights across series
& Solve one scalar problem per distinct $(n,d,s)$ & Reuse $w$ for each new $y,\mu$ & Proposed timed comparison\\
\bottomrule
\end{tabular}
\end{table*}

\subsection{What can and cannot be verified from the current repository}
The implementation \texttt{src/trendestimation/smoothness\_weights.py}
already evaluates direct log-odds weights and reconstructs the trend
without forming $\lambda$. Its tests compare results against the
existing \texttt{lambda\_for\_smoothness()+fit()} path for
difference orders $1$ through $6$, fixed nonzero $\mu$,
the observed-difference plug-in and endpoint cases.
These tests support finite-precision agreement on their sampled cases,
not a globally valid numerical error bound.
No new numerical experiments are performed in this manuscript,
and no numerical tables in the preceding two CSDA drafts are
reclassified as a direct-weight runtime comparison.
